\subsection{最大元与最小元}

如果存在一个元素 \(a\) 属于有序集 \(\mathbf{E}\) ，使得 \(a \leqslant x\) 对于所有 \(x \in \mathbf{E}\) 成立，那么 \(a\) 是 \(\mathbf{E}\) 中具有此性质的唯一元素；因为如果还有 \(b \leqslant x\) 对于所有 \(x \in \mathbf{E}\) 成立，那么我们就应有 \(a \leqslant b\) 与 \(b \leqslant a\) ，从而 \(a = b\) .

定义 4. 设 E 是一个有序集。元素 \(a \in E\) 被称为 E 的最小（相应地，最大）元素，如果对所有 \(x \in E\) 我们有 \(a \leqslant x\) （相应地 \(x \leqslant a\) ）.

一个有序集不一定有最大元素，也不一定有最小元素。如果 E 有最小元素 a，则 a 是关于相反序的最大元素。

如果 E 有最小元素 \(a\) ，那么 \(a\) 是 E 的唯一极小元素；因为如果 \(x \in E\) 不同于 \(a\) ，我们就有 \(a < x\) .

示例

\begin{enumerate}
\def\labelenumi{(\arabic{enumi})}
\item
  设 \(\mathfrak{S}\) 为集合 \(\mathfrak{P}(\mathbf{E})\) （集合 \(\mathbf{E}\) 的子集所成的集合）的一个非空子集。如果 \(\mathfrak{S}\) 关于包含关系有最小（相应地，最大）元素A，则A是 \(\mathfrak{S}\) 中这些集合的交集（相应地，并集）。反之，若 \(\mathfrak{S}\) 中这些集合的交集（相应地，并集）属于 \(\mathfrak{S}\) ，则它是 \(\mathfrak{S}\) 的最小（相应地，最大）元素。
\item
  特别是， \(\varnothing\) 是最小元素，而E是最大元素，均在 \(\mathfrak{P}(E)\) 中。在集合 \(\Phi(E, F)\) （从E的子集到F的映射所成的集合）中，按映射的延拓排序（第1号，例3），空映射是最小元素，且除非F由单一元素组成，否则不存在最大元素。对角线 \(\Delta\) ，即 \(E \times E\) 的对角线，是 E 上等价关系的图所成集合（或 E 上的预序所成集合）的最小元素。
\end{enumerate}

命题 3. 设 E 是一个有序集，并设 E' 是 E 与一个集合 \{a\} （由单个元素组成）的不交并。那么存在唯一的E'上的序，它诱导出E上给定的序，并且使得a是E'的最大元素。

事实上，如果 G 是 E 上序关系的图，那么 \(E'\) 上满足命题条件的序关系之图，必为并集 \(G'\)，即 G 与所有二元组 \((x, a)\) 所成集合的并，其中 \(x \in E'\)；反之，显然 \(G'\) 是 \(E'\) 上一个满足给定条件的序关系之图。¶ 称有序集 \(E'\) 是由 E 添加一个最大元 a 得到的（参见练习 3）。

预序集 E 的子集 A 称为 E 的共尾子集（相应地，共始子集），如果对每个 \(x \in E\) ，存在 \(y \in A\) 使得 \(x \leqslant y\) （相应地 \(y \leqslant x\) ）。说有序集 E 有最大元（相应地，最小元）因此意味着 E 有一个由单个元素组成的共尾（相应地，共始）子集。

